% Zumdahl Chapter 11 -- Properties of Solutions. ELABORATED notes.
% 3 examinable blocks (11.1-11.3) + 1 clearly marked enrichment block
% (11.4-11.7 colligative properties), which are off the current CED.
\documentclass{apchem-notes}

\apcunitword{Chapter}
\apcunit{11}
\apcunittitle{Properties of Solutions}
\apcdoctitle{Guided Notes}
\apcsubtitle{Zumdahl \S11.1--11.3 (CED) $+$ \S11.4--11.7 (enrichment) \textbullet\ PDF pp.~551--581 \textbullet\ 4 blocks}

\begin{document}
\apctitleblock

\begin{apcnote}[Read this first --- about half the chapter is off-syllabus]
\textbf{Examinable: \S11.1--11.3.} Solution composition, the energetics of
dissolving, and the factors that affect solubility --- CED 3.7 and 3.10,
taught in Unit 3.

\textbf{Not examinable: \S11.4--11.8.} Vapour pressure lowering,
boiling-point elevation, freezing-point depression, osmotic pressure, and
colloids. Collectively the \emph{colligative properties}, all dropped from
the AP framework in the 2019 redesign. They appear here as
\textbf{one clearly marked enrichment block}, included because freezing-point
depression is the reason salt is spread on icy roads and osmosis is the
reason cells burst in pure water --- worth an hour, worth zero exam points.

Note that \textbf{molality} exists almost entirely to serve colligative
calculations. It is defined in Block 1 for completeness and because
distinguishing it from molarity sharpens what molarity actually means.
\end{apcnote}

% =====================================================================
\blocklesson{Solution Composition}{Zumdahl \S11.1}

\begin{objectives}
  \item Compute molarity, molality, mass percent, and mole fraction.
  \item Explain why molarity is temperature-dependent and molality is not.
\end{objectives}

\reading{Zumdahl \S11.1}{552--555}

\begin{warmup}[3]
  \item Moles in \SI{58.44}{\gram} \ce{NaCl}:
        \answerblank[2.6cm]{\SI{1.000}{\mole}}
  \item $[\ce{Cl-}]$ in \SI{0.20}{\Molar} \ce{CaCl2}:
        \answerblank[2.6cm]{\SI{0.40}{\Molar}}
  \item Mole fraction of A if $n_A = 0.25$, $n_{\text{tot}} = 1.00$:
        \answerblank[2cm]{0.25}
\end{warmup}

\segment{INSTRUCTION A}{25}{Four ways to state a concentration}

\notesectionz{The definitions}{11.1}
\practice{5}

\renewcommand{\arraystretch}{1.6}
\noindent\begin{tabular}{@{}p{3.0cm}p{5.4cm}p{5.0cm}@{}}
\toprule
\textbf{Measure} & \textbf{Definition} & \textbf{Note}\\
\midrule
\term{Molarity} ($M$) &
  \answerblank[4.2cm]{$\dfrac{\text{mol solute}}{\text{L solution}}$} &
  the working unit of AP chemistry\\
\term{Molality} ($m$) &
  \answerblank[4.4cm]{$\dfrac{\text{mol solute}}{\text{kg solvent}}$} &
  used for colligative work\\
\term{Mass percent} &
  \answerblank[4.6cm]{$\dfrac{\text{mass solute}}{\text{mass solution}}\times100\%$} &
  common in industry and medicine\\
\term{Mole fraction} ($\chi$) &
  \answerblank[4.2cm]{$\dfrac{n_A}{n_{\text{total}}}$} &
  used for gas partial pressures\\
\bottomrule
\end{tabular}

\begin{apctrap}
Two denominators, two different things. Molarity uses litres of
\emph{solution}; molality uses kilograms of \emph{solvent}. Dissolving a
solid changes the volume but not the mass of solvent --- which is exactly
why these are different numbers, and why you cannot convert between them
without knowing the density.
\end{apctrap}

\notesub{Why molality is temperature-independent}

Volume \answerblank[2cm]{expands} when a solution is warmed, so molarity
\answerblank[2cm]{decreases} even though nothing was added or removed. Mass
does not change with temperature, so molality stays
\answerblank[1.8cm]{constant}. That is the only reason molality exists.

\begin{workedexample}[Worked example: one solution, four numbers]
Dissolve \SI{5.00}{\gram} of \ce{NaCl} ($M = 58.44$) in \SI{100.0}{\gram} of
water. The resulting solution has a volume of \SI{102}{\milli\liter}.
\begin{align*}
  n(\ce{NaCl}) &= \frac{5.00}{58.44} = \SI{0.0856}{\mole}\\
  \text{molarity} &= \frac{0.0856}{0.102} = \SI{0.839}{\Molar}\\
  \text{molality} &= \frac{0.0856}{0.1000} = \SI{0.856}{\mole\per\kilo\gram}\\
  \text{mass \%} &= \frac{5.00}{105.0}\times100 = 4.76\%\\
  n(\ce{H2O}) &= \frac{100.0}{18.02} = \SI{5.549}{\mole},\quad
  \chi(\ce{NaCl}) = \frac{0.0856}{5.635} = 0.0152
\end{align*}
Four correct descriptions of the same bottle. Molarity and molality are
close here only because the solution is dilute and aqueous --- do not assume
that in general.
\end{workedexample}

\segment{GUIDED PRACTICE}{15}{Composition drill}

A solution contains \SI{12.0}{\gram} of \ce{NaOH} ($M = 40.00$) dissolved in
\SI{250.0}{\gram} of water; the solution volume is \SI{255}{\milli\liter}.

\begin{enumerate}[leftmargin=*,itemsep=0.35em]
  \item Moles of \ce{NaOH}: \answerblank[2.8cm]{\SI{0.300}{\mole}}
  \item Molarity: \answerblank[2.8cm]{\SI{1.18}{\Molar}}
  \item Molality: \answerblank[2.8cm]{\SI{1.20}{\mole\per\kilo\gram}}
  \item Mass percent: \answerblank[2.8cm]{4.58\%}
\end{enumerate}

\segment{INSTRUCTION B}{20}{Choosing the right measure}

\notesectionz{Which one, and when}{11.1}
\practice{4}

\begin{itemize}[leftmargin=*,itemsep=0.25em]
  \item \textbf{Molarity} for anything involving reaction stoichiometry ---
        you measure out a \answerblank[1.6cm]{volume} and need moles.
  \item \textbf{Mole fraction} for gas mixtures --- it converts directly to
        \answerblank[3.2cm]{partial pressure} via
        $P_A = \chi_A P_{\text{total}}$.
  \item \textbf{Mass percent} when the solute's identity or molar mass is
        unknown, or irrelevant.
  \item \textbf{Molality} essentially only for colligative properties.
\end{itemize}

\segment{APPLICATION}{20}{Composition reasoning}

\begin{enumerate}[leftmargin=*,itemsep=0.4em]
  \item A solution is prepared and then warmed from \SI{20}{\celsius} to
        \SI{60}{\celsius}. State what happens to its molarity, its molality,
        and its mass percent. \answerlines[3]{Molarity decreases --- the
        volume expands while the moles stay fixed. Molality is unchanged ---
        both moles and solvent mass are unaffected by temperature. Mass
        percent is unchanged for the same reason.}
  \item A student prepares a solution by adding \SI{0.100}{\mole} of solute
        to exactly \SI{1.000}{\liter} of water and labels it
        \SI{0.100}{\Molar}. Identify the error and state whether the true
        molarity is higher or lower. \answerlines[3]{Molarity requires
        litres of \emph{solution}, not solvent. Adding solute makes the
        total volume slightly more than 1.000 L, so the true molarity is
        slightly \emph{lower} than 0.100 M. (The molality, however, is
        exactly 0.100.)}
\end{enumerate}

\begin{exitticket}
Why can you not convert molarity to molality without knowing the solution's
density? \answerlines[2]{Molarity is per litre of solution and molality is
per kilogram of solvent; density is what relates the solution's volume to
its mass, so without it the two denominators cannot be connected.}
\end{exitticket}

% =====================================================================
\blocklesson{The Energetics of Dissolving}{Zumdahl \S11.2}

\begin{objectives}
  \item Break dissolving into three energy steps and predict the sign of
        each.
  \item Use $\Delta H_{\text{soln}}$ to explain why some things dissolve and
        others do not.
\end{objectives}

\reading{Zumdahl \S11.2}{555--559}

\begin{warmup}[3]
  \item Molality of \SI{0.50}{\mole} in \SI{500}{\gram} solvent:
        \answerblank[3cm]{\SI{1.0}{\mole\per\kilo\gram}}
  \item Which measure is temperature-independent?
        \answerblank[2.4cm]{molality}
  \item Strongest IMF in water: \answerblank[3.8cm]{hydrogen bonding}
\end{warmup}

\segment{INSTRUCTION A}{25}{Three steps}

\notesectionz{Zumdahl's model of solution formation}{11.2}
\practice{5}

Dissolving is treated as three imaginary steps. They are not a mechanism ---
they are an energy accounting scheme, exactly like Hess's law.

\renewcommand{\arraystretch}{1.45}
\noindent\begin{tabular}{@{}p{1.2cm}p{5.6cm}p{6.4cm}@{}}
\toprule
\textbf{Step} & \textbf{What happens} & \textbf{Sign of $\Delta H$}\\
\midrule
1 & Separate the \answerblank[1.8cm]{solute} into its components &
  \answerblank[2.7cm]{endothermic} --- forces must be overcome\\
2 & Expand the \answerblank[1.8cm]{solvent} to make room &
  \answerblank[2.7cm]{endothermic} --- same reason\\
3 & Let solute and solvent \answerblank[1.6cm]{interact} &
  \answerblank[2.2cm]{exothermic} --- new attractions form\\
\bottomrule
\end{tabular}

\[
  \Delta H_{\text{soln}} = \answerblank[3.6cm]{$\Delta H_1 + \Delta H_2 + \Delta H_3$}
\]

The overall sign can go either way, and small differences between large
numbers decide it.

\begin{apcnote}
For an ionic solute the accounting has familiar names: Step 1 is the
\term{lattice energy} (large and endothermic to overcome), and Step 3 is the
\term{hydration energy} (large and exothermic). Dissolving happens when
hydration nearly pays for the lattice. \ce{NaCl} is a near-tie:
$\Delta H_{\text{soln}} = \SI{+3.9}{\kilo\joule\per\mole}$, barely
endothermic --- which is why table salt dissolves readily but the water gets
very slightly cooler.
\end{apcnote}

\segment{GUIDED PRACTICE}{15}{Predict the thermal effect}

\begin{enumerate}[leftmargin=*,itemsep=0.35em]
  \item \ce{NH4NO3} has $\Delta H_{\text{soln}} =
        \SI{+25.7}{\kilo\joule\per\mole}$. What do you feel when it
        dissolves? \answerblank[4.5cm]{the container gets cold}
  \item \ce{CaCl2} has $\Delta H_{\text{soln}} =
        \SI{-82.8}{\kilo\joule\per\mole}$. What do you feel?
        \answerblank[4.5cm]{the container gets hot}
  \item Which of the two would you put in an instant cold pack?
        \answerblank[3cm]{\ce{NH4NO3}}
\end{enumerate}

\segment{INSTRUCTION B}{20}{Why oil and water do not mix}

\notesectionz{The energetic explanation of ``like dissolves like''}{11.2}
\practice{6}

\begin{workedexample}[Worked example: Zumdahl's oil slick]
Why is oil insoluble in water? Work through the three steps.

\textbf{Step 1} (expand the oil): oil molecules are held only by dispersion
forces --- but they are large molecules, so this is
\emph{moderately} endothermic.

\textbf{Step 2} (expand the water): this requires breaking water's
\answerblank[3.4cm]{hydrogen bonds}, which is \emph{strongly} endothermic
--- the largest term in the accounting.

\textbf{Step 3} (let them interact): a polar water molecule can induce a
dipole in a nonpolar oil molecule, so this is exothermic --- but a
dipole--induced dipole attraction is much
\answerblank[1.6cm]{weaker} than the hydrogen bonds that were sacrificed in
Step 2.

The exothermic Step 3 cannot pay for the endothermic Step 2, so
$\Delta H_{\text{soln}}$ is strongly positive and the two liquids stay
separate.
\end{workedexample}

\begin{apcnote}
This is what ``like dissolves like'' actually \emph{means}. It is not a
rule to memorize --- it is the observation that Step 3 can only repay Step 2
when the new solute--solvent attractions are of the same kind and strength
as the solvent--solvent attractions being broken.
\end{apcnote}

\segment{APPLICATION}{20}{Energetic reasoning}

\begin{enumerate}[leftmargin=*,itemsep=0.4em]
  \item Explain, using the three-step model, why \ce{NaCl} dissolves in
        water but not in hexane. \answerlines[3]{In both cases Step 1
        requires overcoming \ce{NaCl}'s very large lattice energy. Water's
        Step 3 produces strong ion--dipole attractions that largely repay
        it. Hexane is nonpolar and can offer only weak dispersion
        interactions with the ions, so Step 3 comes nowhere near repaying
        Step 1.}
  \item Ethanol is miscible with water in all proportions. Explain in terms
        of the three steps. \answerlines[3]{Step 2 breaks water's hydrogen
        bonds and Step 1 breaks ethanol's, both endothermic. But Step 3 lets
        ethanol's O--H hydrogen bond directly to water --- the same type and
        strength of attraction that was broken --- so the energy is
        essentially repaid and mixing proceeds freely.}
\end{enumerate}

\begin{exitticket}
A solute dissolves even though $\Delta H_{\text{soln}}$ is positive. What
does that tell you about the process? \answerlines[2]{Enthalpy is not the
whole story --- the increase in disorder on mixing can drive dissolving even
when it costs energy. (Unit 9 makes this precise with entropy.)}
\end{exitticket}

% =====================================================================
\blocklesson{Factors Affecting Solubility}{Zumdahl \S11.3}

\begin{objectives}
  \item Predict solubility from molecular structure.
  \item State how pressure and temperature affect solid and gas solubility.
\end{objectives}

\reading{Zumdahl \S11.3}{559--563}

\begin{warmup}[3]
  \item Sign of $\Delta H$ for Step 2 (expanding solvent):
        \answerblank[2.6cm]{positive}
  \item Cold pack salt: \answerblank[2.6cm]{\ce{NH4NO3}}
  \item Which dissolves \ce{I2} better, water or \ce{CCl4}?
        \answerblank[2.4cm]{\ce{CCl4}}
\end{warmup}

\segment{INSTRUCTION A}{25}{Structure effects}

\notesectionz{Hydrophobic and hydrophilic}{11.3}
\practice{6}

\begin{workedexample}[Worked example: two vitamins]
Zumdahl contrasts vitamins A and C.

\textbf{Vitamin A} is built almost entirely from carbon and hydrogen, whose
electronegativities are similar --- so the molecule is essentially
\answerblank[1.8cm]{nonpolar}, or \term{hydrophobic}. It dissolves in body
fat but not in water.

\textbf{Vitamin C} carries many polar O--H and C--O bonds, making it
\answerblank[1.6cm]{polar}, or \term{hydrophilic}, and water-soluble.

The consequence is medical: fat-soluble vitamins (A, D, E, K)
\answerblank[2.4cm]{accumulate} in fatty tissue, so excess intake can cause
hypervitaminosis. Water-soluble vitamins are
\answerblank[1.8cm]{excreted} and must be consumed regularly --- which is
why sailors developed scurvy without fresh food.
\end{workedexample}

\notesub{The chain-length effect}

For alcohols, water solubility \answerblank[1.6cm]{falls} as the carbon
chain grows: methanol and ethanol are miscible, but 1-octanol is nearly
insoluble. The polar --OH group stays the same size while the nonpolar
hydrocarbon portion grows, so the molecule becomes progressively more
hydrophobic overall.

\segment{GUIDED PRACTICE}{15}{Predict solubility}

More soluble in water? Give the structural reason.
\begin{enumerate}[leftmargin=*,itemsep=0.3em]
  \item \ce{CH3OH} or \ce{C6H14}:
        \answerblank[6.4cm]{\ce{CH3OH} --- H-bonds with water}
  \item \ce{NaCl} or \ce{I2}:
        \answerblank[5.5cm]{\ce{NaCl} --- ion--dipole}
  \item \ce{C2H5OH} or \ce{C8H17OH}:
        \answerblank[7.2cm]{\ce{C2H5OH} --- shorter nonpolar chain}
\end{enumerate}

\segment{INSTRUCTION B}{20}{Pressure and temperature}

\notesectionz{What changes solubility}{11.3}
\practice{4}

\renewcommand{\arraystretch}{1.4}
\noindent\begin{tabular}{@{}p{2.4cm}p{4.8cm}p{5.8cm}@{}}
\toprule
 & \textbf{Solids in water} & \textbf{Gases in water}\\
\midrule
Pressure & essentially \answerblank[1.9cm]{no effect} &
  solubility \answerblank[1.8cm]{increases} with pressure\\
Temperature & \emph{usually} increases, but
  \answerblank[2.2cm]{not always} &
  solubility \answerblank[1.8cm]{decreases} with temperature\\
\bottomrule
\end{tabular}

\begin{workedexample}[Worked example: why soda goes flat]
A sealed bottle is pressurized with \ce{CO2} at \SI{5.0}{\atm}. Using
Zumdahl's Henry's law constant $k = \SI{3.1e-2}{\mole\per\liter\per\atm}$:
\begin{align*}
  \text{sealed:}\quad C &= (3.1\times10^{-2})(5.0) = \SI{0.16}{\mole\per\liter}\\
  \text{opened:}\quad C &= (3.1\times10^{-2})(4.0\times10^{-4}) = \SI{1.2e-5}{\mole\per\liter}
\end{align*}
Once opened, the \ce{CO2} equilibrates with the atmosphere, where its
partial pressure is only $4.0\times10^{-4}$ atm. The dissolved concentration
falls by a factor of more than ten thousand --- which is precisely what
``going flat'' means.
\end{workedexample}

\begin{apcnote}
AP expects the \emph{relationships} --- gas solubility rises with pressure
and falls with temperature --- not the Henry's law constant calculation. The
qualitative reasoning is what gets tested: warm soda goes flat faster, and
warm river water holds less dissolved oxygen, which is why thermal pollution
harms fish.
\end{apcnote}

\begin{apctrap}
``Solids always dissolve better when hot'' is false, and Zumdahl says so
explicitly. Most solids do, but sodium sulfate and cerium sulfate become
\emph{less} soluble as temperature rises. Also distinguish two different
claims: heating always makes dissolving \emph{faster}, but it does not
always increase \emph{how much} will dissolve.
\end{apctrap}

\segment{APPLICATION}{20}{Solubility reasoning}

\begin{enumerate}[leftmargin=*,itemsep=0.4em]
  \item Explain why a warm carbonated drink fizzes more violently when
        opened than a cold one. \answerlines[2]{Gas solubility decreases
        with temperature, so the warm liquid holds less dissolved \ce{CO2}
        at equilibrium --- more gas escapes, and faster, when the pressure
        is released.}
  \item A power plant returns warmed water to a river. Explain the effect on
        dissolved oxygen and why it matters.
        \answerlines[3]{Gas solubility falls as temperature rises, so the
        warmed water holds less dissolved \ce{O2}. Fish and aquatic
        organisms depend on that oxygen, so thermal pollution can suffocate
        them even though nothing chemical was added.}
  \item A student says ``heating always dissolves more solute.'' Give a
        precise correction. \answerlines[2]{Heating always makes dissolving
        occur faster, but the maximum amount that dissolves may increase or
        decrease depending on the substance --- sodium sulfate is a
        counterexample.}
\end{enumerate}

\begin{exitticket}
Divers breathing compressed air at depth must ascend slowly. Connect this to
gas solubility. \answerlines[2]{At high pressure more nitrogen dissolves in
the blood. Ascending too quickly drops the pressure rapidly, and the
nitrogen comes out of solution as bubbles --- decompression sickness.}
\end{exitticket}

% =====================================================================
\blocklesson{ENRICHMENT: Colligative Properties}{Zumdahl \S11.4--11.7 --- NOT on the CED}

\begin{apcnote}[This block is optional]
Nothing here is assessed on the AP exam --- colligative properties were
dropped in the 2019 redesign. It is included because these effects explain
several everyday phenomena, and because the central idea (a property that
depends only on how many particles are present, not what they are) is
elegant. Skip it entirely without penalty.
\end{apcnote}

\begin{objectives}
  \item Define a colligative property.
  \item Describe vapour pressure lowering, boiling-point elevation,
        freezing-point depression, and osmotic pressure.
\end{objectives}

\reading{Zumdahl \S11.4--11.7}{563--579}

\segment{INSTRUCTION A}{25}{Counting particles, not identifying them}

\notesectionz{What colligative means}{11.4}
\practice{1}

A \term{colligative} property depends only on the
\answerblank[5.4cm]{number of solute particles} in solution --- not on their
chemical identity. One mole of dissolved sugar and one mole of dissolved
neon would behave identically.

\renewcommand{\arraystretch}{1.4}
\noindent\begin{tabular}{@{}p{4.2cm}p{4.0cm}p{5.0cm}@{}}
\toprule
\textbf{Property} & \textbf{Effect of solute} & \textbf{Everyday example}\\
\midrule
Vapour pressure & lowered & ---\\
Boiling point & \answerblank[1.6cm]{raised} & salted cooking water\\
Freezing point & \answerblank[1.6cm]{lowered} & salt on icy roads;
  antifreeze\\
Osmotic pressure & develops across a membrane & cells swell in pure water\\
\bottomrule
\end{tabular}

\notesub{Why electrolytes count double (or triple)}

Because these properties count \emph{particles}, a mole of \ce{NaCl}
produces \answerblank[1.2cm]{2} moles of particles and a mole of \ce{CaCl2}
produces \answerblank[1.2cm]{3}. That is why salt is far more effective per
gram at melting ice than sugar would be.

\segment{GUIDED PRACTICE}{15}{Particle counting}

Rank by freezing-point depression, greatest first, at equal molality:
glucose, \ce{NaCl}, \ce{CaCl2}.
\answerblank[8.5cm]{\ce{CaCl2} $>$ \ce{NaCl} $>$ glucose (3, 2, 1 particles)}

\segment{INSTRUCTION B}{20}{Osmosis}

\notesectionz{Water moving through a membrane}{11.6}
\practice{6}

A \term{semipermeable} membrane lets solvent through but blocks solute.
Water moves from the \answerblank[2.4cm]{more dilute} side toward the
\answerblank[3.9cm]{more concentrated} side. The pressure needed to stop
that flow is the \term{osmotic pressure}.

\begin{apcnote}
Consequences worth knowing: a red blood cell placed in pure water swells and
bursts, because water floods in toward the higher solute concentration
inside. Placed in concentrated brine it shrivels, as water leaves. This is
why intravenous fluids must be \emph{isotonic} with blood --- and why
saltwater is undrinkable.
\end{apcnote}

\segment{APPLICATION}{20}{Explaining familiar effects}

\begin{enumerate}[leftmargin=*,itemsep=0.4em]
  \item Explain why spreading salt on an icy road melts the ice.
        \answerlines[2]{Dissolved ions lower the freezing point of water
        below the ambient temperature, so ice that was stable at that
        temperature melts.}
  \item Explain why a cucumber shrivels into a pickle in brine.
        \answerlines[2]{The brine has a far higher solute concentration than
        the cucumber's cells, so water flows out of the cells by osmosis.}
\end{enumerate}

\begin{exitticket}
Why is \ce{CaCl2} preferred over \ce{NaCl} for de-icing at very low
temperatures? \answerlines[2]{It releases three particles per formula unit
instead of two, producing a greater freezing-point depression per mole
dissolved.}
\end{exitticket}

\end{document}
